% Generated by mceg equations; edit the contract, then regenerate.
% Requires amsmath. Authored math is retained; TeX is not executed here.
% Definitions and hashes do not establish scientific validity.

% equation-sha256: b0d9f50ebb0751fdb5038d3ec262188f578a8ed1a8c1f6da8bced6299a64ebe3
\paragraph{Equation state\_and\_eos (revision 1)}
The bound fractions and proton fraction conserve hydrogen nuclei and charge.
\begin{equation}
\begin{aligned}
\sum_i f_i+x=1,\quad n_i=f_i n_{\rm H},\quad n_e=n_p=xn_{\rm H},\\\rho=m_{\rm H}n_{\rm H},\quad P=(1+x)n_{\rm H}k_{\rm B}T,\quad \mu=(1+x)^{-1}
\end{aligned}
\label{eq:escape.coupled_5f_hydrogen:state_5f_and_5f_eos:1}
\end{equation}
\noindent\textbf{Scope:} Pure atomic hydrogen; H-minus, when enabled, is a trace opacity approximation and does not alter charge/EOS.

\begin{description}
\item[$f_i$] Bound population divided by total nuclear density (1; ID: fractions).
\item[$x$] Ionized fraction (1; ID: ionized).
\item[$n_{\rm H}$] Total hydrogen-nucleus density (cm\textasciicircum{}-3; ID: density).
\item[$T$] Gas temperature (K; ID: temperature).
\item[$P$] Gas pressure (erg cm\textasciicircum{}-3; ID: pressure).
\end{description}

\noindent\textbf{Assumptions:}
\begin{itemize}
\item Electron mass is neglected. No H2, helium or metal charge donors.
\item The H-minus trace approximation requires a separately checked small abundance.
\end{itemize}

% equation-sha256: add38451d2303aaccc579b1b4b7ab3a543c8024558109379a976feddee32088b
\paragraph{Equation material\_energy (revision 1)}
Total material energy includes excitation and ionization, with ground-state neutral H as the reference.
\begin{equation}
\begin{aligned}
\epsilon=\sum_i f_i E_i+x\chi_{\rm H},\quad a^2=\frac{P}{\rho},\\e=\frac32(1+x)k_{\rm B}T+\epsilon,\quad h=\frac52a^2+\frac{\epsilon}{m_{\rm H}},\quad \mathcal B=\frac{v^2}{2}+h+\Phi
\end{aligned}
\label{eq:escape.coupled_5f_hydrogen:material_5f_energy:1}
\end{equation}
\noindent\textbf{Scope:} Translational monatomic particles plus the represented atomic populations.

\begin{description}
\item[$E_i$] Excitation energy of bound state i above ground (erg; ID: excitation).
\item[$\chi_{\rm H}$] Ground-state ionization energy (erg; ID: binding).
\item[$e$] Internal energy per hydrogen nucleus (erg; ID: energy).
\item[$h$] Specific material enthalpy (erg g\textasciicircum{}-1; ID: enthalpy).
\item[$\mathcal B$] Specific total material energy flux divided by mass flux (erg g\textasciicircum{}-1; ID: bernoulli).
\end{description}

\noindent\textbf{Assumptions:}
\begin{itemize}
\item Collisional excitation and ionization are internal energy redistribution, not additional external cooling terms.
\end{itemize}

% equation-sha256: 9d6db83ed52771790fee3bb3e8a5bdc8bb1e7d97d765064b84a3cb15e9ff586d
\paragraph{Equation atomic\_balance (revision 1)}
The same rate network supplies local equilibrium and population-advection source terms.
\begin{equation}
\begin{aligned}
Q_i=\sum_{j\ne i}(n_jP_{ji}-n_iP_{ij})+n_en_p\alpha_i+n_e^2n_pq_i^{(3)}-n_i(\Gamma_i+n_eq_i^{\rm ion})
\end{aligned}
\label{eq:escape.coupled_5f_hydrogen:atomic_5f_balance:1}
\end{equation}
\noindent\textbf{Scope:} Resolved lower nl states and the explicitly chosen collapsed upper shells.

\begin{description}
\item[$Q_i$] Net creation rate of bound state i (cm\textasciicircum{}-3 s\textasciicircum{}-1; ID: production).
\item[$P_{ij}$] Bound-bound rate from i to j (s\textasciicircum{}-1; ID: rate).
\item[$\alpha_i$] Radiative recombination coefficient including stimulated emission (cm\textasciicircum{}3 s\textasciicircum{}-1; ID: alpha).
\item[$q_i^{(3)}$] Three-body recombination coefficient (cm\textasciicircum{}6 s\textasciicircum{}-1; ID: q3).
\end{description}

\noindent\textbf{Assumptions:}
\begin{itemize}
\item No omitted-level cascade or Case-B correction is added implicitly.
\item The one-way two-photon channel prevents a blanket full-atom Planck-LTE claim.
\end{itemize}

% equation-sha256: 7050da097a5be8bb1dd725f15273f16db7d3e431f51cb6d0e0b0a579d747cc39
\paragraph{Equation radiation\_averages (revision 1)}
Complete redistribution uses the same line profile for radiative rates and transfer coefficients.
\begin{equation}
\begin{aligned}
J_\nu=\frac12\int_{-1}^{1}I_\nu\,d\zeta,\quad H_\nu=\frac12\int_{-1}^{1}\zeta I_\nu\,d\zeta,\quad K_\nu=J_\nu/3,\\\overline J_{ul}=\int\phi_{ul,\nu}J_\nu\,d\nu,\quad P_{lu}^{\rm rad}=B_{lu}\overline J_{ul},\quad P_{ul}^{\rm rad}=A_{ul}+B_{ul}\overline J_{ul},\\B_{ul}=\frac{c^2A_{ul}}{2h\nu_{ul}^3},\qquad g_lB_{lu}=g_uB_{ul}
\end{aligned}
\label{eq:escape.coupled_5f_hydrogen:radiation_5f_averages:1}
\end{equation}
\noindent\textbf{Scope:} Static radiation, isotropic emission and Eddington angular closure; no velocity-induced frequency shifts.

\begin{description}
\item[$J_\nu$] Mean specific intensity (erg s\textasciicircum{}-1 cm\textasciicircum{}-2 sr\textasciicircum{}-1 Hz\textasciicircum{}-1; ID: mean).
\item[$H_\nu$] First intensity moment; flux Fnu=4pi Hnu (erg s\textasciicircum{}-1 cm\textasciicircum{}-2 sr\textasciicircum{}-1 Hz\textasciicircum{}-1; ID: flux\_moment).
\item[$\phi_{ul,\nu}$] Line profile normalized over frequency (Hz\textasciicircum{}-1; ID: profile).
\end{description}

\noindent\textbf{Assumptions:}
\begin{itemize}
\item Eddington closure is an approximation; exact angular or beam irradiation is not established by conservation.
\item The stellar boundary is specified as an incoming intensity, not silently identified with the mean intensity.
\item Einstein B here multiplies mean specific intensity, not radiation energy density; the factor c\textasciicircum{}2/(2h nu\textasciicircum{}3) fixes the convention.
\end{itemize}

% equation-sha256: 6ed0934803af7747858dff1e4b8a604f5fb89c7076ee86a9e4c08570ae0c618c
\paragraph{Equation radiation\_moments (revision 1)}
Spherical Eddington moments use the same radial geometry as the material-flow equations.
\begin{equation}
\begin{aligned}
\frac{1}{r^2}\frac{d(r^2H_\nu)}{dr}=j_\nu-\chi_\nu J_\nu,\qquad \frac{dJ_\nu}{dr}=-3\chi_\nu H_\nu
\end{aligned}
\label{eq:escape.coupled_5f_hydrogen:radiation_5f_moments:1}
\end{equation}
\noindent\textbf{Scope:} Stationary transfer without Doppler shifts, aberration or radiation-pressure force.

\begin{description}
\item[$\chi_\nu$] Net extinction including stimulated emission (cm\textasciicircum{}-1; ID: opacity).
\item[$j_\nu$] Isotropic emissivity (erg s\textasciicircum{}-1 cm\textasciicircum{}-3 sr\textasciicircum{}-1 Hz\textasciicircum{}-1; ID: emission).
\item[$r$] Radial coordinate (cm; ID: radius).
\end{description}

\noindent\textbf{Assumptions:}
\begin{itemize}
\item For a plane-parallel validation problem replace r-squared area by one.
\item Zero extinction must be handled as a transport limit, not by an arbitrary opacity floor.
\end{itemize}

% equation-sha256: 000514c02263e319e155b48ed3e9635decdd2936ec9c61214328f82f9ae12ea5
\paragraph{Equation radiation\_boundaries (revision 1)}
Two incoming characteristics close the finite radiation boundary-value problem.
\begin{equation}
\begin{aligned}
J_\nu(r_0)+\sqrt3H_\nu(r_0)=I_{\nu,\rm in}^{+},\qquad J_\nu(r_1)-\sqrt3H_\nu(r_1)=I_{\nu,\rm in}^{-}
\end{aligned}
\label{eq:escape.coupled_5f_hydrogen:radiation_5f_boundaries:1}
\end{equation}
\noindent\textbf{Scope:} Outgoing intensities are solved; incoming spectra and the finite outer radius are inputs.

\begin{description}
\item[$I_{\nu,\rm in}^{+}$] Incoming outward intensity at the inner boundary (erg s\textasciicircum{}-1 cm\textasciicircum{}-2 sr\textasciicircum{}-1 Hz\textasciicircum{}-1; ID: inner).
\item[$I_{\nu,\rm in}^{-}$] Incoming inward intensity at the outer boundary (erg s\textasciicircum{}-1 cm\textasciicircum{}-2 sr\textasciicircum{}-1 Hz\textasciicircum{}-1; ID: outer).
\end{description}

\noindent\textbf{Assumptions:}
\begin{itemize}
\item A one-direction input Iminus=2 W Bnu gives mean W Bnu in a transparent slab with zero opposite input.
\item There is no independent outgoing-boundary condition.
\end{itemize}

% equation-sha256: e348b0cd081962e803e6afe75f6cf7adb63a57b0b5231e6dd360c4548639639b
\paragraph{Equation radiation\_finite\_volume (revision 1)}
The zeroth-moment equation conserves radiative energy over each finite shell.
\begin{equation}
\begin{aligned}
A_{k+1/2}H_{\nu,k+1/2}-A_{k-1/2}H_{\nu,k-1/2}=V_k(j_{\nu,k}-\chi_{\nu,k}J_{\nu,k}),\\A_{k\pm1/2}=r_{k\pm1/2}^2,\quad V_k=\frac{r_{k+1/2}^3-r_{k-1/2}^3}{3}
\end{aligned}
\label{eq:escape.coupled_5f_hydrogen:radiation_5f_finite_5f_volume:1}
\end{equation}
\noindent\textbf{Scope:} J and process coefficients are cell quantities; H lives at edges; geometric factors exclude 4pi.

\begin{description}
\item[$A$] Radial area factor, or one in a slab (cm\textasciicircum{}2 in spherical geometry; ID: area).
\item[$V$] Shell-volume factor, or cell width in a slab (cm\textasciicircum{}3 in spherical geometry; ID: volume).
\end{description}

\noindent\textbf{Assumptions:}
\begin{itemize}
\item The first moment is discretized with consistent half-cell optical depths and boundary characteristics.
\item A cell-mean rate field, not an unrelated midpoint approximation, supplies the heating budget.
\end{itemize}

% equation-sha256: 989ce1dc636e2a4e09d52ee7e96a7f7fa372ebecceab9a70a60d11660a7815c6
\paragraph{Equation radiative\_heating (revision 1)}
Net radiative heating is consistent with the total-material-energy convention.
\begin{equation}
\begin{aligned}
\mathcal H_{\rm RT}=4\pi\int(\chi_\nu J_\nu-j_\nu)\,d\nu,\quad C_{2\gamma}=n_{2s}A_{2\gamma}(E_{2s}-E_{1s}),\\\mathcal H=\mathcal H_{\rm RT}-C_{2\gamma}
\end{aligned}
\label{eq:escape.coupled_5f_hydrogen:radiative_5f_heating:1}
\end{equation}
\noindent\textbf{Scope:} Two-photon decay is an explicitly escaping channel outside the sampled transfer grid.

\begin{description}
\item[$\mathcal H$] External gain of material energy per unit volume (erg cm\textasciicircum{}-3 s\textasciicircum{}-1; ID: heating).
\item[$C_{2\gamma}$] Escaping two-photon radiative cooling (erg cm\textasciicircum{}-3 s\textasciicircum{}-1; ID: escape).
\end{description}

\noindent\textbf{Assumptions:}
\begin{itemize}
\item Do not subtract collisional energy terms again.
\item If two-photon emissivity is later included in jnu, remove the separate escape subtraction.
\item H-minus and free-free contributions must use the same extinction/emission in transfer and heating.
\end{itemize}

% equation-sha256: c738e8aae822592349124e8d777d74eab951c8144c9caa87e8eabf68cdfd0342
\paragraph{Equation continuum\_closures (revision 1)}
Free-free and optional H-minus continuum are explicit closures, not a hidden chemistry solver.
\begin{equation}
\begin{aligned}
j_\nu^{\rm ff}=\chi_\nu^{\rm ff}B_\nu(T),\qquad n_{\rm H^-}=S_-(T)n_{1s}n_e,\qquad j_\nu^{-}=\chi_\nu^{-}B_\nu(T)
\end{aligned}
\label{eq:escape.coupled_5f_hydrogen:continuum_5f_closures:1}
\end{equation}
\noindent\textbf{Scope:} Maxwellian free-free with the chosen Gaunt approximation; H-minus in a conditional LTE trace closure.

\begin{description}
\item[$n_{\rm H^-}$] Trace H-minus density used by the continuum approximation (cm\textasciicircum{}-3; ID: negative).
\item[$S_-(T)$] LTE attachment equilibrium factor (cm\textasciicircum{}3; ID: saha\_minus).
\end{description}

\noindent\textbf{Assumptions:}
\begin{itemize}
\item H-minus is off by default until its small-abundance and detailed process bookkeeping are checked.
\item No dynamic H-minus chemistry or H-minus feedback on charge and internal energy is claimed.
\end{itemize}

% equation-sha256: 1d6715d545362362d48272c4a221e584c7b55a7df18007766a7e6436afe42256
\paragraph{Equation mass\_and\_species\_flow (revision 1)}
Spherical continuity and material advection replace an imposed local equilibrium closure.
\begin{equation}
\begin{aligned}
\dot M=4\pi r^2m_{\rm H}n_{\rm H}v=\mathrm{constant},\qquad v\frac{df_i}{dr}=\frac{Q_i}{n_{\rm H}}
\end{aligned}
\label{eq:escape.coupled_5f_hydrogen:mass_5f_and_5f_species_5f_flow:1}
\end{equation}
\noindent\textbf{Scope:} Steady one-dimensional radial flow; the local-SE limit is a separately checked limiting regime.

\begin{description}
\item[$\dot M$] Mass flux fixed by finite-branch inflow data (g s\textasciicircum{}-1; ID: mass\_loss).
\item[$v$] Outward material speed (cm s\textasciicircum{}-1; ID: velocity).
\end{description}

\noindent\textbf{Assumptions:}
\begin{itemize}
\item Local equilibrium Qi=0 is not imposed simultaneously on an arbitrary advecting population profile.
\end{itemize}

% equation-sha256: 643bbd77757f6ba06063b8180df8023b3852f402bb7c1c680830dbba1fa1993e
\paragraph{Equation momentum\_and\_energy (revision 1)}
Momentum and total material-energy flux retain composition-dependent pressure and enthalpy.
\begin{equation}
\begin{aligned}
v\frac{dv}{dr}=-\frac1\rho\frac{dP}{dr}-\frac{d\Phi}{dr},\qquad\frac{d\mathcal B}{dr}=\frac{\mathcal H}{\rho v},\\\Phi=-\frac{GM}{r}-\frac{\kappa r^2}{2}
\end{aligned}
\label{eq:escape.coupled_5f_hydrogen:momentum_5f_and_5f_energy:1}
\end{equation}
\noindent\textbf{Scope:} A finite radial branch; optional linear tides, no Roche geometry and no radiation-pressure force.

\begin{description}
\item[$\Phi$] Gravitational plus linear-tidal potential (erg g\textasciicircum{}-1; ID: potential).
\item[$\kappa$] Prescribed outward linear-tidal acceleration coefficient (s\textasciicircum{}-2; ID: tide).
\item[$\mathcal B$] Specific material energy flux (erg g\textasciicircum{}-1; ID: bernoulli).
\end{description}

\noindent\textbf{Assumptions:}
\begin{itemize}
\item Variable mu follows ionization; its derivatives are part of the same residual.
\end{itemize}

% equation-sha256: 3a98f21b03d63bb081a8179d691a175692e47e92d580253090fe88db066e6555
\paragraph{Equation flow\_finite\_volume (revision 1)}
A first-order finite branch uses shared shell volumes for reaction and energy exchanges.
\begin{equation}
\begin{aligned}
f_{i,R}-f_{i,L}=\frac{V}{r_R^2}\left(\frac{Q_i}{n_{\rm H}v}\right)_R,\qquad \mathcal B_R-\mathcal B_L=\frac{V}{r_R^2}\left(\frac{\mathcal H}{\rho v}\right)_R,\\\left[v+\frac{a^2}{v}\right]_L^R=\frac{a_R^2(r_R^2-r_L^2)-GM(r_R-r_L)+\kappa(r_R^4-r_L^4)/4}{r_R^2v_R}
\end{aligned}
\label{eq:escape.coupled_5f_hydrogen:flow_5f_finite_5f_volume:1}
\end{equation}
\noindent\textbf{Scope:} Right-endpoint material state represents each radiation cell; refine the grid before drawing accuracy conclusions.

\begin{description}
\item[$V$] Same shell-volume factor used by radiation (cm\textasciicircum{}3; ID: volume).
\item[$\Delta r$] Physical radial cell width (cm; ID: cell\_width).
\end{description}

\noindent\textbf{Assumptions:}
\begin{itemize}
\item The effective species/energy increment V/rR-squared gives exact discrete total-flux exchange with the radiation equation.
\item The momentum increment integrates pressure and gravitational/tidal geometry with the material coefficient held at the right endpoint.
\end{itemize}

% equation-sha256: eda0e6aeebc5b33c849445269a28613edfe87932ce7fc3c0d3a84372f9052381
\paragraph{Equation finite\_inflow\_and\_sonic\_limit (revision 1)}
Finite inflow data close the initial branch problem; sonic regularity belongs to a distinct transonic problem.
\begin{equation}
\begin{aligned}
n_{\rm H}(r_0)=n_0,\quad T(r_0)=T_0,\quad v(r_0)=v_0,\quad f_i(r_0)=f_{i,0},\\v_c^2=\gamma a_c^2,\qquad \frac{2\gamma a_c^2}{r_c}-\Phi_c^\prime-(\gamma-1)\left(\frac{\mathcal H}{\rho v}-\frac{\epsilon^\prime}{m_{\rm H}}\right)_c=0,\quad\gamma=5/3
\end{aligned}
\label{eq:escape.coupled_5f_hydrogen:finite_5f_inflow_5f_and_5f_sonic_5f_limit:1}
\end{equation}
\noindent\textbf{Scope:} The initial implementation must detect or reject a critical crossing; it does not select mass loss from a sonic boundary.

\begin{description}
\item[$f_{i,0}$] Positive normalized inflow fractions (1; ID: inflow).
\item[$r_c$] Unknown sonic radius in a future transonic BVP (cm; ID: sonic).
\end{description}

\noindent\textbf{Assumptions:}
\begin{itemize}
\item For the future transonic problem, mass flux/base speed and the critical radius are unknowns; do not also prescribe every inflow quantity and sonic regularity.
\end{itemize}

% equation-sha256: de8335080417610a5969e55b83a9857ab9eab7fbf6cac8d63b8f41375b8aba69
\paragraph{Equation joint\_residual\_and\_derivative (revision 1)}
A coupled solution and its derivative must use the same complete residual.
\begin{equation}
\begin{aligned}
\mathcal F(y;\theta)=\begin{pmatrix}R_{\rm RT}\\R_{\rm populations}\\R_{\rm momentum}\\R_{\rm energy}\\R_{\rm boundaries}\end{pmatrix}=0,\qquad \mathcal F_y\,\delta y=-\mathcal F_\theta\,\delta\theta
\end{aligned}
\label{eq:escape.coupled_5f_hydrogen:joint_5f_residual_5f_and_5f_derivative:1}
\end{equation}
\noindent\textbf{Scope:} Only include the equations of the selected closed problem; static SE, thermal equilibrium and advecting flow are distinct configurations.

\begin{description}
\item[$y$] Independent material and radiation unknowns (mixed cgs/log coordinates; ID: state).
\item[$\theta$] Boundary, atomic and geometric model inputs (parameter dependent; ID: parameters).
\end{description}

\noindent\textbf{Assumptions:}
\begin{itemize}
\item A frozen-radiation population Jacobian is not the derivative of the coupled atmosphere.
\item Convergence, local Jacobian regularity and independent directional checks are required before inference.
\end{itemize}

% equation-sha256: 2e4b1987091f0be8f62ebe7574e175b30e9077551551d36a5029a1207fd48d08
\paragraph{Equation total\_energy\_gate (revision 1)}
Total material plus transported and explicitly escaping radiation energy provides a coupled acceptance check.
\begin{equation}
\begin{aligned}
\left[\dot M\mathcal B+4\pi r^2F_{\rm rad}\right]_{r_0}^{r_1}+4\pi\int_{r_0}^{r_1}r^2C_{2\gamma}\,dr=0,\qquad F_{\rm rad}=4\pi\int H_\nu\,d\nu
\end{aligned}
\label{eq:escape.coupled_5f_hydrogen:total_5f_energy_5f_gate:1}
\end{equation}
\noindent\textbf{Scope:} The finite-volume form uses exactly the same frequency weights, source terms and shell volumes.

\begin{description}
\item[$F_{\rm rad}$] Frequency-integrated radiative flux (erg s\textasciicircum{}-1 cm\textasciicircum{}-2; ID: flux).
\end{description}

\noindent\textbf{Assumptions:}
\begin{itemize}
\item This conservation check does not establish the accuracy of angular closure, atomic data or astrophysical boundary inputs.
\end{itemize}

% equation-sha256: 4bbebb03cef98eda154cd8c9b02a0687ced7508e478b7260329fc3729fe759bc
\paragraph{Equation line\_transfer\_coefficients (revision 1)}
Individual transition centers and both lifetimes define static CRD Voigt profiles.
\begin{equation}
\begin{aligned}
\chi_\nu^{\rm bb}=\sum_{ul}\frac{c^2A_{ul}}{8\pi\nu_{ul}^2}\left(\frac{g_u}{g_l}n_l-n_u\right)\phi_{ul,\nu},\qquad j_\nu^{\rm bb}=\sum_{ul}\frac{h\nu_{ul}A_{ul}}{4\pi}n_u\phi_{ul,\nu},\\\int\phi_{ul,\nu}\,d\nu=1,\quad\sigma_{\nu,ul}=\frac{\nu_{ul}}c\sqrt{\frac{k_{\rm B}T}{m_{\rm H}}},\quad\gamma_{\nu,ul}=\frac{A_u^{\rm tot}+A_l^{\rm tot}}{4\pi}
\end{aligned}
\label{eq:escape.coupled_5f_hydrogen:line_5f_transfer_5f_coefficients:1}
\end{equation}
\noindent\textbf{Scope:} The common grid normalizes each finite-window profile; its raw integral remains a resolution diagnostic.

\begin{description}
\item[$T$] Gas temperature (K; ID: temperature).
\item[$n_i$] Resolved/group bound populations (cm\textasciicircum{}-3; ID: bound).
\item[$\phi_{ul,\nu}$] Common opacity, emission and rate profile (Hz\textasciicircum{}-1; ID: profile).
\item[$\mathcal A$] Declared finite atomic arrays (mixed cgs; ID: atom).
\end{description}

\noindent\textbf{Assumptions:}
\begin{itemize}
\item Unlike the earlier observer-ray model, new internally coupled transfer retains the separate branch frequencies.
\item No Stark or pressure broadening, bulk Doppler transport or partial redistribution.
\end{itemize}

% equation-sha256: 8c0592e82287b00618721358f1136cb0313d029244e76607f93b3776709e8c86
\paragraph{Equation bound\_free\_transfer\_coefficients (revision 1)}
The same Kramers/Milne kernel supplies bound-free rates, extinction and emission.
\begin{equation}
\begin{aligned}
\sigma_{i\nu}=\frac{2.815\times10^{29}}{\mathfrak n_i^5\nu^3}\,\mathbf1_{h\nu\ge\chi_i},\quad b^*_{i\nu}=\frac{g_i}{2}\left(\frac{h^2}{2\pi m_e k_{\rm B}T}\right)^{3/2}e^{-(h\nu-\chi_i)/(k_{\rm B}T)},\\\chi_\nu^{\rm bf}=\sum_i\sigma_{i\nu}(n_i-n_en_pb^*_{i\nu}),\qquad j_\nu^{\rm bf}=\frac{2h\nu^3}{c^2}\sum_i\sigma_{i\nu}n_en_pb^*_{i\nu}
\end{aligned}
\label{eq:escape.coupled_5f_hydrogen:bound_5f_free_5f_transfer_5f_coefficients:1}
\end{equation}
\noindent\textbf{Scope:} Principal quantum numbers and level densities are distinct; Gaunt factor one and no orbital-sublevel dependence.

\begin{description}
\item[$\mathcal A$] Declared finite atomic arrays (mixed cgs; ID: atom).
\item[$T$] Gas temperature (K; ID: temperature).
\item[$n_i$] Bound population density (cm\textasciicircum{}-3; ID: bound).
\item[$n_e$] Electron number density (cm\textasciicircum{}-3; ID: electron).
\item[$n_p$] Proton number density (cm\textasciicircum{}-3; ID: proton).
\item[$\nu$] Ordinary photon frequency (Hz; ID: frequency).
\item[$\mathfrak n_i$] Principal quantum number of state i (1; ID: principal).
\end{description}

\noindent\textbf{Assumptions:}
\begin{itemize}
\item Stimulated recombination contributes signed extinction; weak gain is preserved within the explicitly checked transport domain, without clipping.
\item Cross sections are approximate; exact inverse identities do not establish atomic-data accuracy.
\end{itemize}

% equation-sha256: 064c062680f510549d6542171c6b11eede1aa2d546e00434f713bcc91a9966dc
\paragraph{Equation higher\_shell\_dipoles (revision 1)}
Hydrogenic dipole integrals and upper statistical weights define added-shell radiative branches.
\begin{equation}
\begin{aligned}
R_{ul}=\int_0^\infty R_{n_ul_u}(r)R_{n_ll_l}(r)r^3\,dr,\qquad A_{ul}=\frac{64\pi^4e^2\nu_{ul}^3}{3hc^3}|R_{ul}|^2\frac{\max(l_u,l_l)}{2l_u+1},\\A_{UL}=\sum_{u\in U,\,l\in L}\frac{g_u}{g_U}A_{ul}
\end{aligned}
\label{eq:escape.coupled_5f_hydrogen:higher_5f_shell_5f_dipoles:1}
\end{equation}
\noindent\textbf{Scope:} Only added n=4,5 shells are collapsed; the original six low nl states retain their source data.

\begin{description}
\item[$n_u$] Upper principal quantum number (1; ID: upper\_n).
\item[$l_u$] Upper orbital quantum number (1; ID: upper\_l).
\item[$n_l$] Lower principal quantum number (1; ID: lower\_n).
\item[$l_l$] Lower orbital quantum number (1; ID: lower\_l).
\item[$\nu_{ul}$] Transition frequency (Hz; ID: frequency).
\item[$A_{ul}$] Spontaneous electric-dipole rate (s\textasciicircum{}-1; ID: rate).
\end{description}

\noindent\textbf{Assumptions:}
\begin{itemize}
\item Only allowed electric-dipole channels are summed; no high-n top-off or Case-B closure.
\item Statistical populations among l sublevels are assumed only inside each added collapsed shell.
\end{itemize}

% equation-sha256: 949331ce925623b2f17c42138c76070986277125344eab2d02aca95595846f62
\paragraph{Equation added\_shell\_inverse\_rates (revision 1)}
Added shell collisions use an explicitly adopted shell approximation and matched inverse processes.
\begin{equation}
\begin{aligned}
q_{u l}=q^{\rm VS80}_{n_u n_l}\frac{A_{u l}}{\sum_{l^\prime:n_{l^\prime}=n_l}A_{u l^\prime}},\quad q_{l u}=\frac{g_u}{g_l}e^{-(E_u-E_l)/(k_{\rm B}T)}q_{u l},\\q_i^{(3)}=q_i^{\rm ion}\frac{g_i}{2}\left(\frac{h^2}{2\pi m_e k_{\rm B}T}\right)^{3/2}e^{\chi_i/(k_{\rm B}T)}
\end{aligned}
\label{eq:escape.coupled_5f_hydrogen:added_5f_shell_5f_inverse_5f_rates:1}
\end{equation}
\noindent\textbf{Scope:} VS80 is the inspected Cloudy C23 evaluation, not an independently verified transcription of the original paper.

\begin{description}
\item[$\mathcal A$] Declared finite atomic arrays (mixed cgs; ID: atom).
\item[$T$] Gas temperature (K; ID: temperature).
\item[$q_{ul}$] De-excitation coefficient (cm\textasciicircum{}3 s\textasciicircum{}-1; ID: rates).
\end{description}

\noindent\textbf{Assumptions:}
\begin{itemize}
\item Original lower-six tabulated collision strengths and l-mixing are retained.
\item Added ionization fits use n=4 (a,b,c)=(1.27e-6,0.053,-10900) and n=5 (5.39e-6,-0.023,-7380) in a*T\textasciicircum{}b*exp(c/T), with T in K.
\item The original VS80 full paper was unavailable; the altered Cloudy adaptation and its notice are retained.
\end{itemize}

% equation-sha256: bb76bb273a4e87a6759210a5b20add2702eb8328599ae1f472b5a8ac15aabfdb
\paragraph{Equation thermal\_equilibrium (revision 1)}
Static thermal equilibrium solves temperature together with rates and internally transported radiation.
\begin{equation}
\begin{aligned}
Q_i=0\quad\text{for every represented bound state},\qquad\mathcal H(T,\{n_i\},J_\nu)=0
\end{aligned}
\label{eq:escape.coupled_5f_hydrogen:thermal_5f_equilibrium:1}
\end{equation}
\noindent\textbf{Scope:} Prescribed nuclear density and incoming intensities; temperature is an unknown and its initial guess is not a boundary condition.

\begin{description}
\item[$n_{\rm H}$] Prescribed nuclear density (cm\textasciicircum{}-3; ID: density).
\item[$T$] Solved temperature (K; ID: temperature).
\item[$\mathcal H$] Full radiative matter-energy input including escaping two-photon loss (erg cm\textasciicircum{}-3 s\textasciicircum{}-1; ID: heating).
\end{description}

\noindent\textbf{Assumptions:}
\begin{itemize}
\item This static equilibrium is not imposed on the finite advecting flow.
\item Gas temperatures must remain within the actual atomic and optional continuum validity ranges.
\end{itemize}

% equation-sha256: 1689a928532ecab0985cc301be16f646845865ea1c2d22e32055bfa62b69115c
\paragraph{Equation radiation\_admissibility (revision 1)}
The initial static moment solver has explicit weak-gain and angular-realizability limits.
\begin{equation}
\begin{aligned}
\max_\nu\sum_k\max(-\chi_{\nu,k},0)\,\Delta r_k\le0.25,\qquad J_{\nu,\mathrm{face}}\ge\sqrt3|H_{\nu,\mathrm{face}}|
\end{aligned}
\label{eq:escape.coupled_5f_hydrogen:radiation_5f_admissibility:1}
\end{equation}
\noindent\textbf{Scope:} A numerical engineering domain for the adopted two-direction closure, not a physical maser threshold.

\begin{description}
\item[$\chi_\nu$] Signed total extinction (cm\textasciicircum{}-1; ID: opacity).
\item[$r$] Increasing cell-edge radii (cm; ID: radius).
\item[$J_\nu$] Reconstructed face mean intensity (erg s\textasciicircum{}-1 cm\textasciicircum{}-2 sr\textasciicircum{}-1 Hz\textasciicircum{}-1; ID: mean).
\item[$H_\nu$] Face first intensity moment (erg s\textasciicircum{}-1 cm\textasciicircum{}-2 sr\textasciicircum{}-1 Hz\textasciicircum{}-1; ID: flux).
\end{description}

\noindent\textbf{Assumptions:}
\begin{itemize}
\item No negative-opacity clipping or artificial positive opacity floor.
\item Nonnegative two-direction intensities are required up to a stated floating-point tolerance.
\item A wide transparent outward-illuminated sphere may fail this angular closure despite an exact linear energy budget.
\end{itemize}
